\section{Data-only diagnostics and a bias-aware interval}\label{sec:diagnostic}
The bias of Section~\ref{sec:anatomy} cannot be seen from the sample's own residuals, because the fit is evaluated where there are no data. We ask whether a score computed from the experiment alone detects it, and whether an interval can account for it. Scores, thresholds and the interval were fixed before the cells that evaluate them were generated, with fresh seeds.

\textbf{Scores.} \emph{D1} (extrapolation check) fits the features on units whose exposure is below its 80th percentile, scales the mean residual of the units above it to the corner, does the same for the lower tail, and scores the absolute difference of the two implied corner errors.
\emph{D2} is half the range of $\that$ over \est{REG-L}, \est{REG-Q} and \est{REG-G}.
\emph{D3} bounds $|h''|$ by $\hat M=2|\hat c_2|+4\hat s_2$ from the coefficient $\hat c_2$ of $g^2$ in \est{REG-Q} and its robust standard error $\hat s_2$, and the bias by $B=\hat M\,\bigl((1-\hat p)^2+\hat p^2\bigr)/2$.
\emph{D4} (degree transport) fits a cubic in $g$, with a treatment term, to the lowest-degree quartile of units only, reads the corner effect off that curve, and scores its absolute difference from \est{REG-L}; it assumes that the response is the same function of exposure at every degree.
Baselines use no outcomes: the standard error implied by the leverage~\eqref{eq:leverage}, and the corner distance $Z$.
The \emph{bias-aware interval} is $\that\pm(1.96\,\widehat{\mathrm{se}}+B)$ with the sandwich standard error; it is valid when the curvature obeys the bound, and a step beyond the sample or an omitted two-hop term is its stated limit.

\textbf{Evaluation.} The 480 cells cross ER, SBM, BA and WS, $d\in\{5,\dots,80\}$, $p\in\{0.1,0.25,0.5,0.75\}$ and six mechanisms (linear, convex, concave, convex with half its weight on the two-hop exposure, threshold, count) at $n=2000$, 120 replicates per cell. In 277 of the 336 cells with resolved bias (82\%) it exceeds $10\%$ of $|\beta|+|\gamma|$, so flagging every cell has precision 0.82.

\textbf{D1--D3 rank cells inside the smooth families and nowhere else.}
Within concave and convex cells, D2 and D3 rank cells by $|\text{bias}|$ with Spearman correlation 0.78--0.80, against 0.58--0.63 for the baselines (Figure~\ref{fig:diag}). Over all 480 cells the correlations are 0.11 (D1), 0.24 (D2), 0.16 (D3), against 0.30 for $Z$; for flagging large bias the areas under the ROC curve are 0.46, 0.61 and 0.56, below the baselines (0.60, 0.67).
They fail where the misspecification is not a smooth curvature: on the threshold response D1 is anti-informative (AUC 0.11), and on the two-hop response all three are at or below chance (AUC 0.35--0.49) while $Z$ reaches 0.77. Within a cell no score tracks the realised error (median Spearman 0.00--0.05).

\textbf{Transporting a low-degree fit ranks the bias where the response does not depend on degree, and only there.}
A second grid of 132 cells on BA, Stanford3, ego-Facebook and an ER control ($p\in\{0.25,0.5\}$, Bernoulli, 100 replicates) evaluates D4, whose predictions were registered before it existed. On the 54 heterogeneous cells whose response does not depend on degree, D4 ranks cells by $|\text{bias}|$ with Spearman correlation 0.82 [0.69, 0.92], against 0.64 for the next best score (D1 on \est{REG-Q}; the intervals overlap, since 54 cells are few) and 0.28 for $Z$; its AUC is 0.87 against 0.66 for $Z$, and its flag has precision 0.85 and recall 0.81 (Figure~\ref{fig:diag}b). As an estimator it is nearly unbiased on the real graphs (median $|\text{bias}|$ 0.002 on Stanford3, 0.009 on ego-Facebook, against 0.108 and 0.093 for \est{REG-L}), but not on BA or ER, where the cubic is extrapolated from a narrow support (RMSE 0.35 against 0.15, and 0.58 against 0.25).
Where the assumption is false, in 54 cells whose indirect effect or count response depends on degree, the correlation is $-0.39$ [$-0.62$, $-0.11$], the flag fires in 50 cells, and the estimate is wrong by 0.4--0.9 in the direction opposite to \est{REG-L}'s. D4 is a diagnostic conditional on an assumption that no experiment can test.

\textbf{The interval is valid inside its class, and wide.}
Coverage is in Appendix Table~\ref{tab:intervals}. For smooth responses inside the curvature rule (linear, concave, convex) the bias-aware interval covers 0.988; for responses that violate it (threshold, two-hop) it covers 0.851, with $B$ below the bias in 19\% of threshold and 21\% of two-hop cells (Figure~\ref{fig:diag}c); for the count response it covers 0.998. The iid, sandwich and bootstrap intervals cover 0.22--0.73, and 0.04--0.09 when bias dominates (241 cells; Appendix Table~\ref{tab:regime}).
The price is width: mean 5.9 with $\gate=2$, 13 times the sandwich and 9 times the iid interval, because $B$ exceeds $|\text{bias}|$ by a median factor 3.2 in bias-dominated cells. Multiplying the iid half-width by one constant (9.2) so that its mean width matches covers 1.00 in all 480 cells: at equal width, plain inflation is at least as good as $B$, which adds only that it is larger where the bias is larger (Spearman 0.78 within concave and convex cells).

\begin{figure}[tb]
\centering
\includegraphics[width=\linewidth]{fig7_diagnostic.pdf}
\caption{Diagnostics (bias of \est{REG-L}). (a) Spearman correlation of each score with $|\text{bias}|$ over the 480 evaluation cells (black) and within each mechanism (colours); bars, 95\% bootstrap over cells; the two lowest scores are baselines that use no outcomes. (b) The same on the 132 cells of the second grid: heterogeneous graphs where the response does not depend on degree (``holds''), where it does (``stress''), and ER (control). (c) The curvature bound $B$ against $|\text{bias}|$ on the 480 cells; points above the line are cells where the bound fails.}
\label{fig:diag}
\end{figure}
